\section{Análisis de Riesgo}
Se identifican los riesgos que podrían afectar la viabilidad del proyecto utilizando una matriz de impacto y probabilidad, complementada con un análisis estático de seguridad y gobernanza (SAST). A continuación, se presenta un desglose de las vulnerabilidades y fallas detectadas tanto a nivel sistémico como por componente.

\subsection{Análisis de Riesgo Integral del Sistema}

\subsubsection{Núcleo Administrativo y Repositorio (Java / PostgreSQL)}
Este componente constituye el centro de persistencia y gestión jurídica del sistema.
\begin{itemize}
    \item \textbf{Compromiso de Datos en Reposo:} El uso de \textit{pgcrypto} (AES-256) es crítico; una gestión deficiente de las llaves maestras en el servidor expondría la identidad de la comunidad politécnica.
    \item \textbf{Fallas en la Cadena de Custodia:} Existe el riesgo de alteración de evidencias multimedia. Aunque se emplean \textit{hashes} SHA-512, vulnerabilidades en el almacenamiento permitirían la sustitución de archivos probatorios.
    \item \textbf{Inyección en Búsquedas Avanzadas:} Riesgo asociado al manejo de datos semi-estructurados (JSONB) y búsquedas de texto completo si las consultas no se parametrizan correctamente.
\end{itemize}

\subsubsection{Microservicio de Inteligencia Artificial (Python)}
Módulo encargado del triaje automático basado en Procesamiento de Lenguaje Natural.
\begin{itemize}
    \item \textbf{Denegación de Servicio (DoS) Semántico:} Los modelos \textit{Sentence-Transformers} son costosos en recursos. Narrativas excesivamente extensas podrían saturar la memoria RAM del contenedor, deteniendo el flujo de admisión.
    \item \textbf{Sesgo en Clasificación Legal:} Un error en la determinación de competencia frente a las exclusiones del Artículo 4º podría rechazar quejas válidas, vulnerando el debido proceso.
\end{itemize}

\subsection{Análisis de Riesgo: Microservicio del Quejoso}
Este componente, desarrollado en Java con Spring Boot, fue evaluado bajo un modelo crítico de negocio. Los resultados del análisis técnico revelan lo siguiente:

\begin{itemize}
    \item \textbf{Indicadores de Estado:} El microservicio presenta un índice de riesgo de 0.51 y un indicador global de 90.54. No obstante, la calificación de seguridad técnica es baja (1 estrella de 5), indicando la necesidad de correcciones urgenteS.
    \item \textbf{Vulnerabilidades Críticas:}
    \begin{itemize}
        \item \textbf{Inyección - Cross-site Scripting (XSS):} Se detectaron 4 fallas de prioridad muy alta (CWE-79) por neutralización inadecuada de entradas en los controladores de autenticación, perfil, quejas y trámites.
        \item \textbf{Debilidad Criptográfica:} Se identificó un defecto de prioridad alta (CWE-330) en la gestión de usuarios debido al uso de generadores de números pseudo-aleatorios que no cumplen con estándares de resistencia ante ataques.
    \end{itemize}
    \item \textbf{Gestión de Componentes:} Se identificaron 73 dependencias externas. Presentan riesgo de obsolescencia medio componentes como \textit{jakarta.inject-api} e \textit{io.jsonwebtoken:jjwt-api}.
    \item \textbf{Referencia:} Los reportes técnicos detallados se encuentran en el \textit{Anexo E: Reportes Técnicos - Microservicio Quejoso}.
\end{itemize}

\subsection{Matriz de Riesgos del Proyecto}
Para priorizar las estrategias de mitigación, se evalúan los riesgos consolidados en la Tabla \ref{tab:matriz_riesgos_tesis}.

\begin{table}[h!]
\centering
\caption{Matriz de Riesgos: Impacto y Probabilidad}
\label{tab:matriz_riesgos_tesis}
\begin{tabular}{|l|c|c|c|}
\hline
\textbf{Descripción del Riesgo} & \textbf{Probabilidad} & \textbf{Impacto} & \textbf{Prioridad} \\ \hline
Inyección XSS en controladores & Alta & Alto & Crítica \\ \hline
Uso de PRNG no seguro & Media & Muy Alto & Alta \\ \hline
Exposición de datos personales & Baja & Muy Alto & Alta \\ \hline
Saturación del motor de IA (DoS) & Media & Medio & Media \\ \hline
Fallo en clasificación Artículo 4º & Media & Muy Alto & Crítica \\ \hline
\end{tabular}
\end{table}

\subsection{Estrategias de Mitigación y Plan de Acción}
\begin{enumerate}
    \item \textbf{Saneamiento de Entradas:} Implementar mecanismos de codificación de salida y validación estricta en todos los controladores para eliminar los riesgos de XSS detectados.
    \item \textbf{Fortalecimiento Criptográfico:} Sustituir los generadores aleatorios estándar por la clase \textit{java.security.SecureRandom} para asegurar la integridad de los tokens y sesiones.
    \item \textbf{Actualización de Dependencias:} Migrar las librerías \textit{jjwt} y \textit{jakarta} a versiones recientes para mitigar vulnerabilidades por obsolescencia.
    \item \textbf{Resiliencia de Infraestructura:} Configurar el escalado automático de los contenedores de IA para manejar cargas variables de procesamiento.
    \item \textbf{Supervisión Humana:} Mantener el modelo \textit{Human-in-the-loop} donde el abogado valide las sugerencias de la IA, asegurando el cumplimiento normativo.
\end{enumerate}